
\section{Computational cost}
\label{app:runtime}
We distinguish basis training, offline preparation, online evolution, and field output.
The tables collect the recorded costs for the case comparisons and their numerical checks.
Times are measured wall times on the RTX 4080 SUPER workstation; each caption specifies
the workload and timer boundary. Superscript R marks synchronized complete-workload measurements with all points and modes. Other entries use the recorded stage timers. Entries marked n.r. have no separate timer, excl. denotes a branch outside the stated comparison, and m.f. denotes matrix-free evaluation. Composite workloads are listed as bundles.
The analytic Fourier basis is available directly; numerical Laplacian eigensolves
belong to offline preparation. Existing Laplacian modes are reused in the
reported controls.

\subsection{Training cost and its components}
Completed 2000-update spectral training costs appear in \docref{Table}{tab:runtime-training} for 2D and \docref{Table}{tab:Pot-training-basis-cost} for 3D. Topology timers cover positive-basis training; the harmonic lift shared within each Potential/Velocity pair is separate. NKB directly predicts velocity without a harmonic training stage; its update, validation, and total timers appear in \docref{Table}{tab:runtime-nkb-training}. The seed-0 Potential-FG run at $K=128$ takes 195.80 minutes; its $K=256$ counterpart and the matched FG, AL, and canonical Potential families are included in the 3D preparation table.

\begin{table}[!htb]\centering\footnotesize\setlength{\tabcolsep}{4pt}
\caption{Two-dimensional spectral training wall time in minutes for complete 2000-update runs. Periodic entries are mean $\pm$ sample SD over three seeds. Topology totals exclude the separately trained harmonic lift shared within each Potential/Velocity pair. Three-dimensional spectral totals are in \docref{Table}{tab:Pot-training-basis-cost}; NKB totals and components are in \docref{Table}{tab:runtime-nkb-training}.}
\label{tab:runtime-training}
\label{tab:k1024-cost}
\begin{tabular}{lrrr}\toprule
Case&$K$&Potential&Velocity\\\midrule
Periodic & 1024 & $106.428\pm3.196$ & $165.760\pm1.636$\\
1 hole & 1024 & 232.61 & 423.97\\
2 holes & 1024 & 235.48 & 376.76\\
3 holes & 1024 & 233.88 & 388.71\\
\bottomrule\end{tabular}\end{table}
\begin{table}[!htb]\centering\footnotesize\setlength{\tabcolsep}{4pt}
\caption{Instrumented 3D Velocity-training components, in minutes, summed over 2000 committed updates and 80 scheduled validations. Fwd./Rev. include analytic derivatives and the reverse pass includes recomputation. The full wall total also includes initialization, progress logging, cleanup, and other uninstrumented overhead; the listed components identify the principal computational stages. Checkpoints are saved every update through $K=256$ and every 25 updates at higher capacities.}
\label{tab:runtime-training-components}

\begin{tabular}{rrrrrrrr}\toprule
$K$&Fwd.&Rev.&Geometry&Ortho.&Validate&Save&Opt./sample\\\midrule
128 & 10.52 & 35.97 & 25.14 & 0.21 & 8.35 & 4.67 & 0.22\\
256 & 20.85 & 72.05 & 25.22 & 0.62 & 10.76 & 9.26 & 0.22\\
512 & 42.17 & 145.88 & 25.61 & 2.20 & 15.95 & 1.82 & 0.35\\
1024 & 89.94 & 298.41 & 29.88 & 8.86 & 28.61 & 3.77 & 0.67\\
\bottomrule\end{tabular}\end{table}
\paragraph{Architecture and derivative workload.}
The 3D velocity implementation assigns an independent width-128, five-hidden-layer sinusoidal representation network (SIREN) {\citep{sitzmann2020implicit}} to each mode: at $K=1024$ it has 70,126,592 trainable parameters. The shared NKB network has 1,188,352 parameters (\docref{Appendix}{app:nkb}). For Velocity, the neural vector potential $A_\theta$ gives $u=\curl A_\theta$ and the velocity-gradient Rayleigh energy require second spatial derivatives of $A_\theta$ and their parameter adjoints. NKB directly predicts $u$ and its losses use first spatial derivatives. At $K=1024$, neural forward and reverse account for 83.3\% of full training wall time; gauge projection is disabled. The cost breakdown connects the derivative workload and network architecture to the measured training time.

\begin{table}[!htb]\centering\footnotesize\setlength{\tabcolsep}{4pt}
\caption{NKB training timers, in seconds. Updates sum the 2000 measured optimization iterations, including sampling, analytic network forward/reverse, losses, Adam, and checkpoint writes within those iterations. Validation sums the 80 scheduled events. Total additionally includes preflight, initialization, and other run overhead. Topology seed is 1 and the 3D seed is 20260823; update totals measure the combined neural forward/reverse and optimization workload.}
\label{tab:runtime-nkb-training}

\begin{tabular}{lrrrrr}\toprule
Case&$K$&Seed&Updates&Validate&Total\\\midrule
Periodic & 1024 & 0 & 151.75 & 3.17 & 159.13\\
Periodic & 1024 & 1 & 151.59 & 3.19 & 158.98\\
Periodic & 1024 & 2 & 151.45 & 3.17 & 158.78\\
1 hole & 1024 & 1 & 153.04 & 3.28 & 160.96\\
2 holes & 1024 & 1 & 154.53 & 3.33 & 162.52\\
3 holes & 1024 & 1 & 159.76 & 3.52 & 167.91\\
3D & 128 & 20260823 & 41.68 & 4.72 & 50.63\\
3D & 256 & 20260823 & 61.90 & 7.01 & 73.30\\
3D & 512 & 20260823 & 121.50 & 13.69 & 139.66\\
3D & 1024 & 20260823 & 325.27 & 37.50 & 367.58\\
\bottomrule\end{tabular}\end{table}

\subsection{Offline basis, projection, and operator preparation}
For bounded domains, basis inference evaluates the frozen NKB model on every assessment node and all modes.
It is the median of three synchronized passes including coordinate/input transfer
and checksum reduction, excluding derivatives and output storage. Projection/preparation
contains the recorded full-mass initial projection and reference/error-moment preparation.
Derived-field preparation contains the derivative caches and linear operators.
Each independently timed pass characterizes its stated workload; production queries reuse the prepared basis and operators.

\begin{table}[!htb]\centering\footnotesize\setlength{\tabcolsep}{4pt}
\caption{Offline costs in seconds. Basis gives all-node field-evaluation times; periodic entries include analytic derivatives. Projection includes the recorded reference preparation where bundled. Derived means field/linear-operator preparation, and $Q$ denotes accepted assembly of the quadratic tensor $\mathcal Q$. The incl. periodic derivative fields are included in basis evaluation, and their nonlinear loads are matrix-free. n.r. marks components without a separate measurement. The $K=512$ derived cost covers all 512 modes, basis-cache transposition, and linear-operator preparation. These columns exclude neural training and online integration.}
\label{tab:runtime-offline}

\begin{tabular}{lrlrrrr}\toprule
Case&$K$&Basis&Basis eval.&Projection&Derived&$Q$\\\midrule
3D & 512 & Velocity & $454.46^{\mathrm R}$ & 18.50 & 3036.39 & 3645.51\\
3D & 1024 & Velocity & $679.90^{\mathrm R}$ & 41.25 & 10381.08 & 32434.41\\
Periodic s0 & 1024 & Potential & 143.21 & 1.73 & \textit{incl.} & \textit{m.f.}\\
Periodic s1 & 1024 & Potential & 142.83 & 1.73 & \textit{incl.} & \textit{m.f.}\\
Periodic s2 & 1024 & Potential & 143.98 & 1.63 & \textit{incl.} & \textit{m.f.}\\
Periodic s0 & 1024 & Velocity & 144.22 & 1.63 & \textit{incl.} & \textit{m.f.}\\
Periodic s1 & 1024 & Velocity & 149.66 & 1.72 & \textit{incl.} & \textit{m.f.}\\
Periodic s2 & 1024 & Velocity & 141.75 & 1.61 & \textit{incl.} & \textit{m.f.}\\
Periodic s0 & 1024 & NKB & 1.36 & 1.68 & \textit{incl.} & \textit{m.f.}\\
Periodic s1 & 1024 & NKB & 1.10 & 1.64 & \textit{incl.} & \textit{m.f.}\\
Periodic s2 & 1024 & NKB & 1.09 & 1.63 & \textit{incl.} & \textit{m.f.}\\
1 hole & 1024 & NKB & 3.55 & 31.34 & $22.58^{\mathrm R}$ & 3247.81\\
2 holes & 1024 & NKB & 14.02 & 118.76 & $90.01^{\mathrm R}$ & 12682.87\\
3 holes & 1024 & NKB & 13.88 & 118.83 & $90.00^{\mathrm R}$ & 12685.85\\
3D & 128 & NKB & 2.99 & 30.18 & 139.57 & 83.94\\
3D & 256 & NKB & 3.50 & 95.56 & 311.72 & 576.10\\
3D & 512 & NKB & 4.67 & 190.53 & 728.16 & 3841.67\\
3D & 1024 & NKB & 7.20 & 322.78 & 1939.69 & 32446.09\\
\bottomrule\end{tabular}\end{table}
For periodic flow, basis evaluation includes the analytic velocity and curl fields on all $256^2$ quadrature nodes; projection includes the full mass/curl-stiffness assembly and generalized eigensolve. The rotational nonlinear load is evaluated matrix-free.
The synchronized 3D Velocity basis pass evaluates analytic curl on all 1,599,754 nodes in the original FP32 spatial arithmetic, with FP64 checksum accumulation; coordinate/input transfer and checksum reduction are included, and storage is excluded. Its checksum agrees with the original full-mass basis energy to $1.10\times10^{-14}$ relative error. The topology basis passes evaluate every field and derivative and assemble the full mass, derivative, and divergence moments on 1,048,576 nodes for one hole and 4,194,304 for two and three holes. Each saved matrix agrees with its original operator to the reported precision.




\subsection{Online evolution, output, and assessment}
\docref{Table}{tab:nkb-runtime} separates coefficient evolution from output costs in 3D.
Laplacian controls use the original grid-based reduced generator, whereas the listed
Velocity and NKB trajectories use assembled coefficient operators. The NS reference
uses its original split viscosity/normal-penalty integrator. The common physical interval and fine time step support comparison of the recorded execution strategies.
Flow errors accompany these timings in the capacity plots. Output measures field
reconstruction, transfer, storage, and recorded diagnostics.

\begin{table}[!htb]\centering\footnotesize\setlength{\tabcolsep}{4pt}
\caption{3D online costs in seconds, $T=1.16$. Coarse/fine steps are $0.00125/0.000625$. ROM output timers cover three selected fields; NS output timers cover its 30 saved fields. Laplacian online timers include the field operations required at every integration step, excluding assessment outside the stepping timer. Repeated NS rows are the measured reference runs for the two capacity experiments; n.a. indicates that NS has no retained-mode count.}
\label{tab:nkb-runtime}

\begin{tabular}{rllrrr}\toprule
$K$&Method&Execution&Coarse&Fine&Output\\\midrule
128 & NKB & Coefficient & 3.57 & 6.84 & $22.30^{\mathrm R}$\\
256 & NKB & Coefficient & 4.90 & 9.78 & $22.26^{\mathrm R}$\\
512 & Velocity & Coefficient & 3.85 & 7.03 & 20.41\\
512 & Laplacian & Grid ROM & 119.45 & 238.71 & $0.26^{\mathrm R}$\\
512 & NKB & Coefficient & 8.29 & 16.58 & 7.98\\
1024 & Velocity & Coefficient & 24.30 & 48.59 & 32.72\\
1024 & Laplacian & Grid ROM & 130.04 & 260.17 & $0.26^{\mathrm R}$\\
1024 & NKB & Coefficient & 766.10 & 1532.83 & 14.56\\
\textit{n.a.} & NS (K512 comparison) & Grid NS & $63.71^{\mathrm R}$ & 124.75 & 7.01\\
\textit{n.a.} & NS (K1024 comparison) & Grid NS & $63.66^{\mathrm R}$ & 126.68 & 7.04\\
\bottomrule\end{tabular}\end{table}
\begin{table}[!htb]\centering\footnotesize\setlength{\tabcolsep}{4pt}
\caption{2D measured execution intervals in seconds, $K=1024$. Each periodic method--seed pair is listed once. Periodic runs reach $T=18\pi$ with 241 saved fields; topology runs reach $T=1$ with 20 initial states. Periodic integration timers include synchronized coefficient capture and exclude field reconstruction and file writing. The matched potential, velocity and NKB runs use $256^2$ operator points and $128^2$ observations. Topology entries time their original 20-state batches.}
\label{tab:runtime-2d-online}

\begin{tabular}{llrll}\toprule
Case&Method&$\Delta t$&Time (s)&Timer scope\\\midrule
Periodic s0 & Potential & $0.005$ & 182.62 & Integration + coefficient capture\\
Periodic s1 & Potential & $0.005$ & 182.63 & Integration + coefficient capture\\
Periodic s2 & Potential & $0.005$ & 182.64 & Integration + coefficient capture\\
Periodic s0 & Velocity & $0.005$ & 185.97 & Integration + coefficient capture\\
Periodic s1 & Velocity & $0.005$ & 183.62 & Integration + coefficient capture\\
Periodic s2 & Velocity & $0.005$ & 181.99 & Integration + coefficient capture\\
Periodic s0 & NKB & $0.005$ & 182.99 & Integration + coefficient capture\\
Periodic s1 & NKB & $0.005$ & 182.81 & Integration + coefficient capture\\
Periodic s2 & NKB & $0.005$ & 182.59 & Integration + coefficient capture\\
1 hole & NKB & $0.001$ & 408.46 & 20-state coefficient batch\\
2 holes & NKB & $0.001$ & 402.47 & 20-state coefficient batch\\
3 holes & NKB & $0.001$ & 418.86 & 20-state coefficient batch\\
\bottomrule\end{tabular}\end{table}
\begin{table}[!htb]\centering\footnotesize\setlength{\tabcolsep}{4pt}
\caption{High-capacity Velocity pipeline wall times, in seconds. Process-stage intervals are grouped by purpose for each 20-state topology experiment. They include process startup, I/O, and process monitoring, and are kept separate from the synchronized NKB coefficient timers in \docref{Table}{tab:runtime-2d-online}. Shared harmonic training is excluded.}
\label{tab:runtime-topology-pipeline}

\begin{tabular}{lrrr}\toprule
Stage&1 hole&2 holes&3 holes\\\midrule
Boundary fields and assessment & 60.05 & 70.07 & 70.04\\
Initial projection & 20.01 & 40.03 & 40.02\\
Linear/operator preparation & 874.27 & 977.54 & 978.13\\
Nonlinear tensor & 3484.71 & 3504.59 & 3501.97\\
ROM execution & 520.06 & 610.08 & 610.08\\
Reference fields and error assessment & 570.13 & 2250.30 & 2280.28\\
\bottomrule\end{tabular}\end{table}
The one-hole kernel comparison in \docref{Appendix}{app:protocol} isolates coefficient stepping from streamed field contraction. Its kernel-level intervals differ from the process-stage workloads in \docref{Table}{tab:runtime-topology-pipeline}.


\subsection{Numerical checks and auxiliary costs}
\begin{table}[!htb]\centering\footnotesize\setlength{\tabcolsep}{4pt}
\caption{Additional 3D measured costs in seconds. Velocity check includes full-grid RHS and one-step comparisons; refined projection uses the matched-step NS reference. Velocity/NKB evaluation bundles include both trajectories, checks, and output and overlap the preceding component timers. NKB capacities time evaluation after operator preparation. Laplacian refresh includes both time steps and fresh FP64 projection. Bundle totals and individual components describe their respective timer boundaries.}
\label{tab:runtime-3d-checks}

\begin{tabular}{rrrrrrr}\toprule
$K$&Velocity load&Velocity check&Refined proj.&Velocity eval.&NKB eval.&Lap. refresh\\\midrule
512 & 0.71 & 183.30 & 19.37 & 386.33 & 69.50 & 376.58\\
1024 & 3.36 & 438.64 & 35.32 & 761.22 & 2337.10 & 407.83\\
\bottomrule\end{tabular}\end{table}
At $K=128,256$, joint Velocity/Laplacian evaluation costs 547.33 and 3449.86 s, covering both methods, both time steps, and their field outputs. The corresponding NKB evaluation bundles take 25.91 and 30.12 s. These totals characterize the complete evaluation workloads at the two capacities.

\paragraph{NKB operator precision.}
The accepted $K=1024$ tensor assembly uses FP64 short-dot products and compensated accumulation. Every online RK4 stage likewise uses compensated products and summation. Their measured costs appear in Tables~\ref{tab:runtime-offline} and~\ref{tab:nkb-runtime}. Independent full-grid RHS comparisons verify the assembled operator to the reported precision (\docref{Table}{tab:nkb-3d-checks}).

\paragraph{Gauge-solver timings.}
\docref{Table}{tab:gauge-solver-ablation} groups setup, forward, reverse, and independent
assessment at the frozen checkpoint, with residuals visualized in \docref{Figure}{fig:gauge-solver-ablation}.
\docref{Table}{tab:gauge-precision} adds source-evaluation times for the precision controls.
The workload evaluates 32 frozen columns with reused geometry-dependent MFS factors;
the listed timers cover the source evaluation, correction, differentiation, and assessment stages.



\subsection{Inference accounting and reuse thresholds}
\label{app:inference-amortization}
Basis-query timing follows the continuous evaluation boundary in \docref{Section}{sec:inference-cost}; numerical construction additionally charges geometry-specific mesh preparation, point mapping, factorization, and the spectral solve. \docref{Table}{tab:runtime-offline} reports all-node evaluation costs, while \docref{Appendix}{app:unseen-geometry} gives fresh numerical constructions on unseen domains. The reuse counts below concern trajectories that share spatial bases and operators.

\begin{table}[!htb]\centering\small
\caption{Reuse thresholds computed from original wall-time receipts, without rounding inputs. A reuse is one $T=1.16$, 1856-step trajectory on the same geometry, with the same time step and existing operators. The baseline is the numerical NS integration timer. Training-only counts set $T_{\rm off}=0$; adding recorded assembly or derived-field preparation charges those stages once. Derived-field costs include the recovered $K=512$ preparation timers. Per-query projection/output and separately unrecorded stages lie outside these integration-only thresholds.}
\label{tab:inference-amortization}
\begin{tabular}{rrrr}\toprule
$K$&Training only&+ $Q$ assembly&+ derived fields\\\midrule
512&123&154&179\\
1024&359&774&907\\
\bottomrule\end{tabular}\end{table}

For a full workload, replace the denominator of \doceqref{eq:inference-break-even} by the difference between complete per-query costs, including initial projection, selected field reconstruction, transfers, and storage. Output costs are assigned according to the requested number of reconstructed fields: the recorded NS workload saves 30 fields, and the ROM workload reconstructs three selected fields.

For basis generation over a geometry family, write $T_{\rm num}^{(g)}$ for a fresh numerical basis construction and $T_{\rm query}^{(g)}$ for the frozen neural map with the same postprocessing. For $G$ geometry queries indexed by $g$, training amortizes when
\begin{equation}
\sum_{g=1}^{G}\left(T_{\rm num}^{(g)}-T_{\rm query}^{(g)}\right)\geq T_{\rm train}.
\end{equation}
If the per-geometry costs are constant, write them as $T_{\rm num}$ and $T_{\rm query}$; then $G_\star=\lceil T_{\rm train}/(T_{\rm num}-T_{\rm query})\rceil$. This geometry-level threshold requires matched basis-generation timers. Each new geometry also incurs its own operator assembly when autonomous dynamics are requested. Arbitrarily dense neural spatial queries preserve the same learned continuous function; their cost scales with requested points and modes, and their accuracy is that of the learned representation.



\paragraph{Tensor-product and nonuniform parameter sweeps.}
A tensor-product sweep with $m$ values for each of $d$ independent query parameters contains $m^d$ cases. The per-axis break-even resolution $r_\star(d)$ and integer sampling count $m_\star(d)$ are
\begin{equation}
r_\star(d)=n_\star^{1/d},\qquad
m_\star(d)=\left\lceil r_\star(d)\right\rceil.
\label{eq:sweep-resolution}
\end{equation}
Here $d$ counts varied physical or query parameters and $n_\star$ is the integration-only reuse count in \doceqref{eq:inference-break-even}.


Unequal axis counts $m_1,\ldots,m_d$ meet the threshold when $\prod_{j=1}^d m_j\geq n_\star$; their geometric mean $(\prod_jm_j)^{1/d}$ is the mean per-axis resolution. Sparse, constrained, or correlated sweeps use the actual query count. With per-query savings $\Delta T_i$ and reusable cost $T_{\rm one\mbox{-}time}$, break-even requires $\sum_i\Delta T_i\geq T_{\rm one\mbox{-}time}$, including per-case preparation. Fixed-geometry initial-condition sweeps reuse the frozen basis and operators; geometry sweeps additionally charge geometry generation and required operator assembly.

